\section{Literature Review}

Prior research in insider threat detection and digital forensics spans four primary domains: graph neural networks for behavioral modeling, file-metadata forensics, provenance-graph cybersecurity, and digital forensic readiness standards.

\textbf{Graph Neural Networks for Insider Threat Detection}: The application of deep graph learning to insider threat detection has gained significant momentum due to the expressive power of graph structures in capturing complex entity interactions. \citet{wei2024ewatcher} introduced E-Watcher, a spatial-temporal graph architecture evaluated on CERT benchmark datasets, reporting a headline accuracy of 98.48\%. Similarly, \citet{li2023tgcnda} developed TGCN-DA, leveraging temporal graph convolutional networks with domain adaptation to achieve an ROC-AUC of 0.9500. \citet{xiao2022mewrgnn} presented MEWRGNN (ROC-AUC 0.9400), using multi-edge weighted relational GNNs to model heterogeneous user activities. More recently, \citet{xiao2024sentinel} proposed SENTINEL (ROC-AUC 0.9300), focusing on dynamic temporal subgraphs. Other baseline approaches, including DTGI \citep{ga02023dtgi} (ROC-AUC 0.9100), Vidhya spatio-temporal modeling \citep{vidhya2024spatiotemporal} (ROC-AUC 0.9100), and sequence-based recurrent neural networks \citep{le2020insider} (ROC-AUC 0.9200, TTD 14.0 min), have established foundational benchmarks. However, existing GNN models predominantly focus on high-volume exfiltration patterns and fail to integrate file-level metadata provenance or court-admissible forensic evidence packaging.

\textbf{File-Metadata Forensics}: A complementary line of research investigates document and image metadata to identify unauthorized file manipulation and identity spoofing. \citet{umair2024exif2vec} introduced Exif2Vec (ROC-AUC 0.8200), embedding image EXIF metadata into vector spaces to detect device manipulation. \citet{yasenenko2025imagemeta} explored image structural headers for anomaly detection (ROC-AUC 0.7600). Forensic analyses of Office Open XML (DOCX) and Portable Document Format (PDF) files by \citet{didriksen2014ooxml} and \citet{fu2015docx} demonstrated that revision identifiers (\texttt{rsid}), document creation timestamps, and author metadata contain critical forensic artifacts capable of exposing ghost-authoring and metadata-stripping attacks. Nevertheless, prior metadata forensic solutions operate on isolated individual files rather than connecting document provenance into a global organizational knowledge graph.

\textbf{Provenance Knowledge Graphs in Cybersecurity}: Knowledge graphs have emerged as a robust paradigm for security context representation. \citet{kurniawan2022krystal} presented KRYSTAL KG (ROC-AUC 0.8900), constructing heterogeneous knowledge graphs from system logs. \citet{hamid2025metashield} introduced MetaShield (ROC-AUC 0.8600) for metadata security monitoring. Earlier rule-based provenance systems, such as PS0 by \citet{mavroeidis2018ps0}, combined W3C-PROV ontology alignment with static rule checking (ROC-AUC 0.7200, F1 0.6100). Standardized provenance models, such as the W3C PROV-O specification \citep{sikos2020prov}, provide formal semantics for tracking entities, activities, and agents, but lack automated integration with deep neural anomaly detectors.

\textbf{Digital Forensic Readiness and ISO/IEC 27043}: Digital Forensic Readiness (DFR) addresses an organization's capability to collect, preserve, and analyze digital evidence before an incident occurs. The ISO/IEC 27043 international standard establishes formal attributes for forensic readiness, including identification, collection, acquisition, preservation, analysis, presentation, chain of custody, and integrity verification \citep{shoderu2025dfrbust}. \citet{shoderu2025dfrbust} introduced DFR-BUST to evaluate readiness matrices in enterprise settings. The cryptographic foundation of tamper-evident log preservation relies on Merkle hash trees and digital signature schemes, as originally formalized by \citet{garfinkel2013forensic}.

\textbf{Summary Gap Statement}: Despite significant advancements across these individual domains, no prior system provides a unified framework that braids cross-format file-metadata harvesting, W3C-PROV knowledge graph representation, boolean rule AST filtering, spatio-temporal GNN anomaly detection, and ISO/IEC 27043-compliant Ed25519-signed Merkle evidence packaging into an end-to-end operational architecture. TURRET OS directly fills this critical research gap.
